% ─────────────────────────────────────────────────────────────────────────────
% Preámbulo LaTeX — Manual de estilo gráfico y editorial de la Dirección Técnica
% (DT 2026, v1.0). Aplica el formato del manual a la salida PDF de la versión
% extendida. Solo formato (tipografía, color, hoja); no toca la redacción.
% ─────────────────────────────────────────────────────────────────────────────

% ── Paquetes heredados del YAML (no eliminar: los anexos usan landscape) ──────
\usepackage{microtype}
\usepackage{etoolbox}
\usepackage{pdflscape}
\apptocmd{\maketitle}{\newpage}{}{}

% ── Paleta institucional (Manual §1.2) ───────────────────────────────────────
\usepackage{xcolor}
\definecolor{tinta}{HTML}{161A1D}      % texto principal y títulos H1/H2
\definecolor{textosec}{HTML}{6B6863}   % texto secundario, notas al pie
\definecolor{gris}{HTML}{98989A}       % gris / marcadores neutros
\definecolor{regla}{HTML}{D8D3C7}      % reglas finas, bordes de tabla
\definecolor{reglasuave}{HTML}{ECE8DE} % rejilla, separación de filas
\definecolor{papel}{HTML}{FBFAF7}      % papel cálido, cajas de cita
\definecolor{guinda}{HTML}{9B2247}     % guinda institucional (H4)
\definecolor{guindaprof}{HTML}{611232} % guinda profundo
\definecolor{verde}{HTML}{1E5B4F}      % verde institucional
\definecolor{verdeprof}{HTML}{002F2A}  % verde profundo
\definecolor{dorado}{HTML}{A57F2C}     % dorado (H3)
\definecolor{arena}{HTML}{E6D194}      % arena

% Color del texto corrido (Normal, Manual §5.1)
\color{tinta}

% ── Jerarquía de títulos (Manual §4) ─────────────────────────────────────────
% H1 20 pt tinta · H2 12 pt tinta · H3 11 pt dorado · H4 9 pt guinda. Todos en
% negrita (toma Noto Sans Bold de mainfont). La negrita la aplican los estilos,
% no se escribe a mano (Manual §9.4).
\usepackage{titlesec}

\titleformat{\section}
  {\bfseries\color{tinta}\fontsize{20}{24}\selectfont}
  {\thesection}{0.6em}{}
\titleformat{\subsection}
  {\bfseries\color{tinta}\fontsize{12}{14}\selectfont}
  {\thesubsection}{0.6em}{}
\titleformat{\subsubsection}
  {\bfseries\color{dorado}\fontsize{11}{14}\selectfont}
  {\thesubsubsection}{0.6em}{}
\titleformat{\paragraph}[runin]
  {\bfseries\color{guinda}\fontsize{9}{12}\selectfont}
  {\theparagraph}{0.6em}{}

% Espaciado vertical razonable alrededor de los títulos
\titlespacing*{\section}{0pt}{18pt}{8pt}
\titlespacing*{\subsection}{0pt}{14pt}{6pt}
\titlespacing*{\subsubsection}{0pt}{12pt}{5pt}
\titlespacing*{\paragraph}{0pt}{10pt}{0.6em}

% ── Notas al pie (Manual §5.3): 8 pt, texto secundario ───────────────────────
% Se redefine solo el texto de la nota para no teñir el \footnotesize que usan
% los anexos en landscape.
\makeatletter
\renewcommand\@makefntext[1]{%
  \parindent 1em\noindent
  \hb@xt@1.8em{\hss\@makefnmark}%
  {\fontsize{8}{13.8}\selectfont\color{textosec}#1}}
\makeatother

% ── Tablas: toque de paleta (tipográfico + paleta) ───────────────────────────
% Reglas de tabla en color regla #D8D3C7 (compatible con booktabs de kable).
\usepackage{colortbl}
\arrayrulecolor{regla}

% ── Citas (Manual §1.2): fondo papel cálido + filete izquierdo guinda ────────
\usepackage{mdframed}
\renewenvironment{quote}
  {\begin{mdframed}[
     backgroundcolor=papel,
     linecolor=guinda,
     linewidth=2pt,
     topline=false, bottomline=false, rightline=false,
     innerleftmargin=12pt, innerrightmargin=12pt,
     innertopmargin=8pt, innerbottommargin=8pt,
     skipabove=8pt, skipbelow=8pt]}
  {\end{mdframed}}
